\documentclass[11pt]{article}
\usepackage[a4paper,margin=29mm]{geometry}
\usepackage[T1]{fontenc}
\usepackage[utf8]{inputenc}
\usepackage{lmodern}
\usepackage{amsmath,amssymb,amsthm}
\usepackage{microtype}
\usepackage[hidelinks]{hyperref}
\hypersetup{pdftitle={A uniform linear bound for facial cycles on surfaces},pdfauthor={Kaiwen Jin}}
\newtheorem{theorem}{Theorem}
\newtheorem{lemma}[theorem]{Lemma}
\newtheorem{proposition}[theorem]{Proposition}
\newtheorem{corollary}[theorem]{Corollary}
\theoremstyle{remark}
\newtheorem{remark}[theorem]{Remark}
\DeclareMathOperator{\girth}{girth}
\title{A uniform linear bound for facial cycles on surfaces}
\author{Kaiwen Jin\\ \small\href{mailto:yui.ui.siki@gmail.com}{yui.ui.siki@gmail.com}}
\date{October 15, 2026}
\begin{document}
\maketitle
\begin{abstract}
Let $G$ be a finite simple graph of girth at least $\ell\ge3$, maximal
in its fixed crossing-free embedding on an orientable surface with
$g\ge1$ handles or a nonorientable surface with $g\ge1$ crosscaps.
We prove that every simple cycle forming the entire boundary of a face
has length at most $100000g\ell$. The same constant works for both
surface types, including noncellular embeddings.
The proof bounds the patterns of equal-radius balls along one facial
cycle by counting alternating disjoint trees. Published hypergraph
covering estimates then provide few centers, whose boundary runs are
short by the girth condition.
\end{abstract}

\section{Introduction}

Fix a graph embedded without crossings on a surface and an integer
$\ell\ge3$. An embedding is maximal subject to girth at least $\ell$
if adding any edge between previously nonadjacent vertices without
crossings creates a cycle of length less than $\ell$. This condition
restricts the geometry of a face even when the number of vertices in
the graph is unrestricted. We consider a simple cycle whose image is
the entire boundary of one face, and call it a facial cycle.

Axenovich, Kie{\ss}le, Sagdeev and Zhukovskii
\cite[Theorem~2]{Axenovich} proved the upper bound
$24(2g+1)\ell^2$ for a surface with $g$ handles or $g$ crosscaps.
P\'alv\"olgyi and Z\'olomy \cite[Theorem~3]{PalvolgyiZolomy} proved that
for each fixed orientable genus the length is at most $c_g\ell$.
The joint-parameter question in \cite[Question~2]{Axenovich} asks
whether an absolute constant bounds the length by $Cg\ell$.
We give such a bound for both types of surface.

\begin{theorem}\label{thm:main}
Let $g\ge1$ and $\ell\ge3$ be integers, and let $\Sigma$ be a connected
closed surface that is either orientable with $g$ handles or
nonorientable with $g$ crosscaps. Let $G$ be a finite simple graph
embedded without crossings on $\Sigma$, of girth at least $\ell$ and
maximal for this fixed embedding subject to that girth constraint.
Every simple cycle $C$ forming the entire boundary of a face satisfies
\[
|C|\le100000g\ell.
\]
\end{theorem}

The graph may be disconnected and the embedding may be noncellular.
The boundary is understood as a point set; the proof also covers the
one-sided case permitted by this interpretation. In particular it
includes the usual case of a simple facial boundary walk.

Saturation makes the graph distance between any two vertices of a
facial cycle at most $\ell-2$. A bound on distance does not yet give a
bound on the number of boundary vertices. Following the ball-piercing
and nearest-center strategy of \cite[Section~4]{PalvolgyiZolomy}, we
cover the boundary by a small number of metric balls and count the
consecutive boundary runs associated with their centers.
The genus dependence comes from an estimate on ball patterns on the
one designated face. Many changes of a fixed site witness would give
two disjoint shortest-path trees meeting that face alternately; Euler's
formula limits the number of such alternations. The resulting cubic
trace count gives a logarithmic VC bound. Together with the fractional
ball-packing theorem of \cite{Bousquet} and the covering bound of
\cite{DSW}, it supplies a number of centers bounded by an absolute
multiple of $g$.

After one subdivision of every edge, the balls have enough radius to
meet pairwise, while the girth makes each induced ball a tree.
Each boundary run is then short. The same topological count bounds
how often the center regions can alternate and yields the theorem.
Sections~\ref{sec:topology} and \ref{sec:balls} give the two counting
estimates and the covering argument; Section~\ref{sec:main-proof}
finishes the proof.

\section{Disjoint subgraphs on one face}\label{sec:topology}

Surfaces are connected and closed unless a boundary is specified.
Graphs are finite, undirected and unweighted. They are simple except for
auxiliary multigraphs explicitly allowed below. A face is a component of
the complement of the embedded graph. A disk face with a simple boundary
means that its closure is a closed disk with that cycle as boundary.
The length of a simple cycle is its
number of edges, which also equals its number of vertices. We write $d_J$
for graph distance in $J$ and
\[
B_J(x,r)=\{y\in V(J):d_J(x,y)\le r\}
\]
for an integer radius $r\ge0$. The Euler genus of a connected closed surface
is $\kappa=2-\chi$, where $\chi$ is its Euler characteristic. It is twice the
number of handles in the orientable case and the number of crosscaps in the
nonorientable case.
The girth is the length of a shortest simple cycle, and is infinite for
an acyclic graph.

We first record the surface calculation used in both counting arguments.
The classification of closed surfaces and Euler's formula apply to both
types of surface; see \cite[Corollary 6.4 and Theorem 6.5]{Erickson}.
Capping the $b$ boundary circles of a compact connected surface $S$ gives a
closed surface $\widehat S$ and
\begin{equation}\label{eq:boundary-euler}
\chi(S)=2-\kappa(\widehat S)-b\le2-b.
\end{equation}
Indeed, each added disk increases the Euler characteristic by one.
Gluing along a circle adds Euler characteristics, since a circle has
Euler characteristic zero.

\begin{lemma}\label{lem:cap}
Let a connected finite graph, possibly with loops and parallel edges, be
embedded on a connected closed surface of Euler
genus $\kappa$. Capping the boundary circles of a connected regular
neighborhood of the graph produces a cellular embedding on a closed
surface of Euler genus at most $\kappa$, with the same local cyclic orders.
\end{lemma}
\begin{proof}
Build the neighborhood $N$ from small vertex disks and thin strips along
edges. Let it have $b$ boundary circles. The complementary surfaces $S_i$
have $b_i\ge1$ boundary circles, with $\sum_i b_i=b$. The capped surface
has Euler genus
\[
\kappa_N=2-\chi(N)-b
 =\kappa+\sum_i\bigl(\chi(S_i)-b_i\bigr)\le\kappa,
\]
by \eqref{eq:boundary-euler}. Each boundary circle follows a facial boundary
walk of the graph. Its collar and the added cap form a disk face.
The vertex disks, and hence their cyclic orders, are unchanged.
\end{proof}

In such a disk-and-strip neighborhood, every edge contributes two sides
to the facial boundary walks. A walk turns between neighboring halfedges
in a vertex disk. This remains true for a twisted strip: with a local
orientation the turn is a successor or a predecessor in the cyclic order.

\begin{lemma}\label{lem:hub}
Suppose that a graph embedded on a closed surface of Euler genus $\kappa$
has a disk face bounded by a simple cycle $D$. Let $q\ge2$ pairwise
vertex-disjoint connected subgraphs contain $R$ distinct selected vertices
of $D$, with at least one selected vertex in each subgraph.
If neighboring selected vertices in circular order have different
subgraph labels, including the last and first, then
\begin{equation}\label{eq:hub}
R\le2q+2\kappa-2.
\end{equation}
In particular, if two disjoint trees meet $D$ at $b$ positive and $b$
negative selected vertices in alternating order, then $b\le\kappa+1$.
\end{lemma}
\begin{proof}
Replace each connected subgraph by a tree joining its selected vertices.
Place a new vertex $h$ in the disk face and draw $R$ disjoint spokes from
$h$ to the selected vertices, in their circular order. Retain only the
trees and spokes, and contract each tree to one vertex. These contractions
are possible inside disk neighborhoods of the trees and do not involve
$h$. We obtain a connected bipartite multigraph $M$ with $q+1$ vertices
and $R$ edges, all incident with $h$. Its cyclic neighbor word at $h$ has
no adjacent equal labels. Parallel spokes are allowed; there are no loops.

Apply Lemma~\ref{lem:cap} to $M$, obtaining a cellular embedding of Euler
genus $\kappa_M\le\kappa$. Every facial boundary walk has positive even
length. A walk of length two would have the form $h,a,h$. If it used two
different spokes, they would have the same label $a$ and would be neighbors
at $h$, contrary to the cyclic word. If it used one spoke twice, the turn
at $h$ would have to return from that halfedge to itself, which requires
$\deg(h)=1$. Here $\deg(h)=R\ge q\ge2$. Thus a doubled bridge also cannot
give a length-two face. The argument concerns corners in the local disk
at $h$, so it also holds on a nonorientable surface.

If there are $f$ faces, it follows that $4f\le2R$. Euler's formula gives
$(q+1)-R+f=2-\kappa_M$, and therefore
$R\le2q+2\kappa_M-2\le2q+2\kappa-2$.
For the last assertion use $q=2$ and $R=2b$.
\end{proof}

\section{Balls centered on the boundary}\label{sec:balls}

The next estimate controls the patterns that equal-radius balls can realize
when their centers lie on one facial cycle. No girth assumption is needed here.

\begin{lemma}\label{lem:traces}
Let $J$ be a connected graph with a disk face bounded by a simple cycle $D$
on a surface of Euler genus $\kappa$. For $m\ge1$ distinct sites
$S\subseteq V(J)$ and an integer $r\ge0$, put
\[
I(v)=\{s\in S:d_J(s,v)\le r\}\qquad(v\in V(D)).
\]
The number of distinct traces $I(v)$ is at most
\begin{equation}\label{eq:traces}
1+m^2\bigl[1+(\kappa+1)(m-1)\bigr].
\end{equation}
For $m\ge2$ this is at most $(\kappa+2)m^3$.
\end{lemma}
\begin{proof}
Choose one boundary representative for each distinct nonempty trace, and
fix an order on $S$. Group the representatives by their trace cardinality
$j$ and the site $a$ farthest from the representative among sites in its
trace, resolving ties by the fixed order. There are at most $m^2$ groups.
Within one group the traces are distinct $j$-sets containing $a$.

A singleton group contributes one representative. In any other group, let
$q$ be the immediate predecessor of a representative $p$ in the group's
circular order. The set $I(p)\setminus I(q)$ is nonempty, since the two
traces are distinct and have equal cardinality. Assign $p$ to one site
$b$ in this difference. Since $a$ belongs to both traces, $b\ne a$.

Fix $a$ and a site $b$ assigned $t\ge1$ times. Denote the corresponding
positive representatives by $p_i$ and their negative predecessors by $q_i$.
The $q_i$ are distinct, because the predecessor operation is a bijection
on the full group. No $q_i$ can equal any $p_j$, since $b$ is absent from
negative traces and present in positive traces. Each pair $q_i,p_i$ is
consecutive within the group, so these $2t$ vertices alternate in their
induced circular order. We have
\begin{equation}\label{eq:positive-negative}
d_J(a,q_i)\le r<d_J(b,q_i),\qquad
d_J(b,p_i)\le d_J(a,p_i)\le r.
\end{equation}

Take the subtree of a fixed breadth-first search tree rooted at $a$
joining all $q_i$, and the corresponding subtree rooted at $b$ joining all
$p_i$. Their root paths are shortest paths in $J$. The two subtrees are
vertex-disjoint. Indeed, if an $a,q_i$-path and a $b,p_j$-path shared $z$,
then
\[
d_J(b,z)+d_J(z,q_i)\ge d_J(b,q_i)>r
 \ge d_J(a,q_i)=d_J(a,z)+d_J(z,q_i).
\]
Hence $d_J(b,z)>d_J(a,z)$, and the positive path would give
\[
d_J(a,p_j)\le d_J(a,z)+d_J(z,p_j)
 <d_J(b,z)+d_J(z,p_j)=d_J(b,p_j),
\]
contradicting \eqref{eq:positive-negative}.
Lemma~\ref{lem:hub} now yields $t\le\kappa+1$.

There are $m-1$ possible sites $b\ne a$. Thus a nonsingleton group has at
most $(\kappa+1)(m-1)$ members, and the bound
$1+(\kappa+1)(m-1)$ also covers singleton groups. Sum over the at most
$m^2$ groups and allow one empty trace to obtain \eqref{eq:traces}.
For $m\ge2$,
\[
1+m^2[1+(\kappa+1)(m-1)]
 =1+(\kappa+1)m^3-\kappa m^2\le(\kappa+2)m^3.
\]
\end{proof}

\subsection{VC dimension}

A finite hypergraph consists of a vertex set $V$ and a set of edges
$\mathcal E\subseteq2^V$. Its VC dimension is the largest cardinality of a
set $S\subseteq V$ for which the intersections $E\cap S$, $E\in\mathcal E$,
give every subset of $S$. Repeated edges are identified.
For the graph and facial cycle of Lemma~\ref{lem:traces}, consider the
hypergraph $\mathcal B_{D,r}$ with vertex set $V(J)$ and edge set
$\{B_J(v,r):v\in V(D)\}$. Its vertices can lie anywhere in $J$.

\begin{corollary}\label{cor:vc}
The VC dimension of $\mathcal B_{D,r}$ satisfies
\begin{equation}\label{eq:vc}
\delta(\mathcal B_{D,r})<16\log_2(\kappa+2).
\end{equation}
\end{corollary}
\begin{proof}
If $m\ge2$ sites are shattered, symmetry of graph distance identifies
$B_J(v,r)\cap S$ with the trace $I(v)$ of Lemma~\ref{lem:traces}. Hence
$2^m\le(\kappa+2)m^3$. Put $a=\log_2(\kappa+2)\ge1$.
At $m=16a$, the ratio of $2^m$ to $(\kappa+2)m^3$ is
$2^{15a}/(16a)^3$. It is $8$ at $a=1$, and its logarithmic derivative
$15\ln2-3/a$ is positive for $a\ge1$. Also $2^m/m^3$ is increasing for
$m\ge16$, since $\ln2-3/m>0$. These inequalities follow already from
$\ln2>2/3$; an explicit verification is
\begin{equation}\label{eq:ln2}
\ln2=\int_1^2\frac{dx}{x}>
\int_1^2\left(\frac43-\frac{4x}{9}\right)dx=\frac23,
\end{equation}
where the difference of the two integrands is $(2x-3)^2/(9x)$.
Thus shattering is impossible for $m\ge16a$. Dimensions zero and one also
satisfy \eqref{eq:vc}. This is the VC dimension of the hypergraph on $V(J)$,
as required for a vertex transversal.
\end{proof}

\subsection{A transversal}

For a hypergraph with nonempty edges, let $\nu$ be the maximum number of
pairwise disjoint edges and let $\tau$ be the minimum cardinality of a set
meeting every edge. Their fractional versions are
\begin{align*}
\nu^*&=\max\left\{\sum_{E\in\mathcal E}w_E:
 w_E\ge0,\ \sum_{E\ni v}w_E\le1\quad(v\in V)\right\},\\
\tau^*&=\min\left\{\sum_{v\in V}x_v:
 x_v\ge0,\ \sum_{v\in E}x_v\ge1\quad(E\in\mathcal E)\right\}.
\end{align*}
Finite linear-program duality gives $\nu^*=\tau^*$; see
\cite[p.~1]{Bousquet}. We write $e=\exp(1)$, and $\ln$ denotes the natural
logarithm throughout.

The following density estimate makes the genus dependence in the
fractional ball bound explicit.
\begin{lemma}\label{lem:density}
Every simple minor of a graph embedded on a closed surface of Euler genus
$\kappa$ has average degree at most $6+\sqrt{6\kappa}$.
\end{lemma}
\begin{proof}
Deleting edges and contracting edges with distinct endpoints preserve
embeddability on the surface; delete loops and parallel duplicates to
obtain a simple minor. Consider a connected minor with $n\ge3$ vertices
and $t$ edges. By Lemma~\ref{lem:cap} it has a cellular embedding of Euler
genus at most $\kappa$. Its facial walks have length at least three.
A walk of length one requires a loop; a length-two walk using different
edges requires parallels. A length-two walk using the same edge requires
degree one at both endpoints, so the connected graph would be a single
edge, contrary to $n\ge3$.

Euler's formula and $3f\le2t$ give $t\le3n-6+3\kappa$.
Simplicity also gives $2t/n\le n-1$. Thus
\[
\frac{2t}{n}\le\min\left\{n-1,\ 6+\frac{6\kappa}{n}\right\}
 \le6+\sqrt{6\kappa}.
\]
For the last inequality, if $n\le\sqrt{6\kappa}+1$ use $n-1$.
Otherwise use the second bound; the case $\kappa=0$ gives $6$ directly.
Components of order one or two have average degree at most one.
For a disconnected minor take the vertex-weighted average of the
component averages.
\end{proof}

We use the following two published hypergraph estimates.

\begin{proposition}[Bousquet et al.]\label{prop:packing}
Let $J$ be a finite simple graph with an edge, and let $d$ be the maximum
average degree of a simple minor of $J$. If $\mathcal F$ has vertex set
$V(J)$ and each of its edges is a ball of $J$, then
\begin{equation}\label{eq:fractional-packing}
\nu^*(\mathcal F)\le e d\,\nu(\mathcal F).
\end{equation}
\end{proposition}
This is the non-edgeless case of \cite[Theorem~4]{Bousquet}. A ball
subfamily is allowed: it need not consist of all balls of $J$.

\begin{proposition}[Ding--Seymour--Winkler]\label{prop:dsw}
Let $\mathcal F$ be a finite hypergraph with at least two distinct nonempty
edges and VC dimension $\delta$. Then
\begin{equation}\label{eq:dsw}
\tau(\mathcal F)\le
2\delta\,\tau^*(\mathcal F)\log_2\bigl(11\tau^*(\mathcal F)\bigr).
\end{equation}
\end{proposition}
We use \cite[(5.4), printed p.~10]{DSW}, with the base-two logarithms
specified in its Section~5. The fractional form comes from integer copies
of vertices of a rational optimal cover. If its value is $\lambda$, choose
$N$ so that all $Nx_v$ are integers, and give $v$ that many copies.
There are $N\lambda$ vertices in total, while every duplicated edge has
at least $N$ vertices. Copies of one vertex cannot be distinguished by
edges, so duplication does not increase VC dimension
\cite[(5.3)]{DSW}. A transversal of the duplicated hypergraph maps to one
of the original. If duplication leaves just one distinct edge, a
transversal of size one suffices; otherwise the uniform bound
\cite[(5.1)]{DSW} applies with $\varepsilon=1/\lambda$.
If the original hypergraph has one nonempty edge, we instead use
$\tau=\tau^*=1$ directly.

\begin{lemma}\label{lem:cover}
Let $g\ge1$ be an integer. Suppose that a connected graph $J$ has a disk
face bounded by a simple cycle $D$ on a closed surface of Euler genus
$\kappa\le2g$. If $r\ge0$ is an integer and
$d_J(u,v)\le2r$ for all $u,v\in V(D)$, then there is a vertex set $Z$
meeting every ball $B_J(v,r)$, $v\in V(D)$, with
\begin{equation}\label{eq:cover}
|Z|<35840g.
\end{equation}
\end{lemma}
\begin{proof}
The midpoint vertex of a shortest path of length at most $2r$ belongs to
both balls centered at its endpoints. Hence the nonempty hypergraph
$\mathcal B_{D,r}$ has $\nu=1$. The graph $J$ contains $D$, so it has an
edge and Proposition~\ref{prop:packing} applies.
By Lemma~\ref{lem:density} and $g\ge1$,
\[
d\le6+\sqrt{6\kappa}\le6+\sqrt{12g}\le10\sqrt g.
\]
Consequently its fractional cover value satisfies
$1\le\lambda=\tau^*=\nu^*\le10e\sqrt g$.

Choose a minimum transversal $Z$ and set $k=|Z|$.
If there is only one distinct edge, $k=1$ proves \eqref{eq:cover}.
Otherwise Proposition~\ref{prop:dsw} and Corollary~\ref{cor:vc} apply.
The function $x\ln(11x)$ is increasing for $x\ge1$, and $\kappa+2\le4g$.
Converting both base-two logarithms gives
\begin{equation}\label{eq:logarithmic-cover}
k\le\frac{320e}{(\ln2)^2}\sqrt g\,\ln(4g)\ln(110e\sqrt g).
\end{equation}

For an explicit uniform bound, put $u=\ln g\ge0$ and $B=\ln(110e)$.
We have $2<B<7$. Indeed, $5/2<e<3$ and $110<(5/2)^6$ give these
inequalities. The coarse bounds on $e$ can be checked by the chord and
tangent bounds for $1/x$:
\[
\ln(5/2)\le\frac5{12}+\frac7{24}+\frac9{40}
 =\frac{14}{15}<1,\qquad
\ln3\ge\frac23+\frac25=\frac{16}{15}>1.
\]
The first uses the chords on $[1,3/2]$, $[3/2,2]$ and $[2,5/2]$;
the second uses tangents at $3/2$ and $5/2$ on $[1,2]$ and $[2,3]$.

We also have $\ln2<1$ by integrating $1/x<1$ on $(1,2]$, so $\ln4<2$.
The derivative
\[
\frac{d}{du}\bigl[(u+\ln4)e^{-u/4}\bigr]
 =e^{-u/4}\left(1-\frac{u+\ln4}{4}\right)
\]
changes from positive to negative at $u=4-\ln4>0$. Hence the maximum
is $4\sqrt2/e$. Also $e^{u/4}\ge1+u/4$ and $B>2$
give $B+u/2\le B e^{u/4}<7e^{u/4}$. Thus \eqref{eq:logarithmic-cover} yields
\[
k<\frac{8960\sqrt2}{(\ln2)^2}g
 <30240g<35840g,
\]
using $\sqrt2<3/2$ and \eqref{eq:ln2}.
\end{proof}

\section{Proof of the main theorem}\label{sec:main-proof}

\begin{proof}[Proof of Theorem~\ref{thm:main}]
Fix an arbitrary facial cycle $C$, and let $F$ be its face.
Any two vertices $x,y$ of $C$ satisfy
\begin{equation}\label{eq:original-distance}
d_G(x,y)\le\ell-2.
\end{equation}
For equal or adjacent vertices this is immediate. Otherwise, join them
by a simple arc through the connected face $F$ and add the missing edge.
Such an arc exists by taking short arcs into incident face sectors and
joining their interior endpoints in $F$. Maximality forces a cycle of
length below $\ell$ through the new edge. Removing that edge gives the
path in \eqref{eq:original-distance}. This is also the boundary-distance
observation of \cite[Lemma~1]{Axenovich}.

Since the entire boundary of $F$ is $C$, the face is open and closed
relative to $\Sigma\setminus C$, and is therefore a whole component of
that complement. If $C$ does not separate $\Sigma$, then
$F=\Sigma\setminus C$ and $G=C$. Equation~\eqref{eq:original-distance}
gives $\lfloor|C|/2\rfloor\le\ell-2$, so $|C|\le2\ell-3$.
This also handles a one-sided curve under the point-set boundary
interpretation: its Moebius-band collar, minus the core curve, is
connected, so the curve is nonseparating.

Now suppose that $C$ separates. Its two sides have compact closures $S$
and $\overline F$, each with one boundary circle $C$, and the interior
of $\overline F$ contains no graph. Replace this side by a disk, capping
$S$. Capping both sides and gluing along $C$ gives
\[
\kappa(\widehat S)+\kappa(\widehat{\overline F})=\kappa(\Sigma).
\]
Both summands are nonnegative. Thus the same graph, with the same graph
distances, is now embedded with a disk face bounded by $C$ on a closed
surface of Euler genus $\kappa\le\kappa(\Sigma)\le2g$.
We retain \eqref{eq:original-distance} for the remainder of the proof.

Subdivide every edge once, and let $H$ be the component containing the
doubled cycle $C^*$. Removing other components changes no relevant
distance or boundary ball. The disk face persists,
$|C^*|=2|C|$, and $\girth(H)\ge2\ell$, since every subdivided simple cycle
projects to an original simple cycle of half its length. Distances between
old vertices are doubled. Every boundary vertex is within one edge of
an old vertex of $C$, so
\begin{equation}\label{eq:subdivided-distance}
d_H(u,v)\le1+2(\ell-2)+1=2\ell-2
 \qquad(u,v\in V(C^*)).
\end{equation}
Put $r=\ell-1$. Then the boundary diameter is at most $2r$, and
$\girth(H)\ge2r+2$. Lemma~\ref{lem:cover} gives a transversal
$Z=\{z_1,\dots,z_k\}$ of its boundary-centered radius-$r$ balls, with
$k<35840g$.

Every induced radius-$r$ ball of $H$ is a tree. It is connected, since
root geodesics stay in the ball. A breadth-first search tree rooted at
its center has depth at most $r$. Any other induced edge, together with
the unique tree path between its endpoints, would create a simple cycle
of length at most $2r+1$, contradicting the girth.

Assign every vertex of $H$ to a nearest $z_j$, breaking ties by the
smallest index. A shortest path from a vertex $v$ to its assigned center
$z_j$ stays in the same region. Indeed, for a vertex $u$ on that path,
if another center were nearer at $u$, or tied there with a smaller index,
the inequality
\[
d_H(v,z_i)\le d_H(v,u)+d_H(u,z_i)
 \le d_H(v,u)+d_H(u,z_j)=d_H(v,z_j)
\]
would give a better choice at $v$, with a strict inequality in the
nearer case. Thus each region contains shortest paths to its center.
For each label used on $C^*$, take the subtree of a fixed rooted
shortest-path tree joining all its assigned boundary vertices.
These trees are pairwise vertex-disjoint.

Because $Z$ meets $B_H(v,r)$ for every $v$ on $C^*$, every boundary
vertex is within $r$ of its assigned center. A proper circular run of
one label is consequently a simple path in the induced ball of that
center. In a tree its length is at most the sum of the distances from
the center to its endpoints, hence at most $2r$ edges or $2r+1$ vertices.
The whole cycle cannot have one label, as that would put $C^*$ in an
induced tree ball.

Let $q\ge2$ be the number of labels used on $C^*$ and let $R$ be the
number of maximal circular runs, including the last/first adjacency.
Choose one boundary vertex from each run. Their neighboring labels
differ, so Lemma~\ref{lem:hub}, applied to the $q$ disjoint trees, gives
$R\le2q+2\kappa-2\le2k+4g-2$. Summing the disjoint run vertex counts,
\[
|C^*|\le(2k+4g-2)(2r+1)
 =(2k+4g-2)(2\ell-1)<71684g(2\ell).
\]
Finally $|C^*|=2|C|$ gives $|C|<71684g\ell\le100000g\ell$, completing
the proof for both surface types and all the stated parameters.
\end{proof}

\section*{AI assistance}
OpenAI Codex was used for literature searches, developing the argument,
drafting, and checking the mathematical text.

\begin{thebibliography}{9}
\bibitem{Axenovich}
M.~Axenovich, L.~Kie{\ss}le, A.~Sagdeev, and M.~Zhukovskii,
\emph{Faces in girth-saturated graphs on surfaces},
Discrete \& Computational Geometry \textbf{76} (2026), 1503--1523.
\href{https://doi.org/10.1007/s00454-026-00833-x}{doi:10.1007/s00454-026-00833-x}.
The theorem and question locators used here refer to
\href{https://arxiv.org/abs/2410.13481v3}{arXiv:2410.13481v3}.

\bibitem{Bousquet}
N.~Bousquet, W.~Cames van Batenburg, L.~Esperet, G.~Joret,
W.~Lochet, C.~Muller, and F.~Pirot,
\emph{Packing and covering balls in graphs excluding a minor},
Combinatorica \textbf{41} (2021), 299--318.
\href{https://doi.org/10.1007/s00493-020-4423-3}{doi:10.1007/s00493-020-4423-3}.
Theorem~4 and page locators refer to
\href{https://arxiv.org/abs/2001.04517v3}{arXiv:2001.04517v3}.

\bibitem{DSW}
G.-L.~Ding, P.~Seymour, and P.~Winkler,
\emph{Bounding the vertex cover number of a hypergraph},
Combinatorica \textbf{14} (1994), 23--34.
\href{https://doi.org/10.1007/BF01305948}{doi:10.1007/BF01305948}.
Statements (5.1), (5.3), and (5.4) and printed page numbers refer to the
\href{https://math.dartmouth.edu/~pw/papers/hyperg.ps}{author-hosted original}.

\bibitem{Erickson}
J.~Erickson, \emph{Surface classification}, lecture~6 of
\emph{Computational Topology}, author course notes, 2009,
distributed in the 2017 course archive.
\url{https://jeffe.cs.illinois.edu/teaching/comptop/2017/notes/surface-classification.pdf}.

\bibitem{PalvolgyiZolomy}
D.~P\'alv\"olgyi and K.~Z\'olomy,
\emph{Stabbing non-piercing sets and face lengths in large girth plane graphs},
2025.
\href{https://arxiv.org/abs/2504.10618v2}{arXiv:2504.10618v2}.
\end{thebibliography}
\end{document}
